\documentclass[10pt,a4paper]{amsart}
\usepackage[margin=19mm]{geometry}
\usepackage{amsmath,amssymb,mathtools}
\usepackage[scr=rsfs,cal=euler]{mathalfa}
\usepackage{xfrac}
\usepackage[unicode,hidelinks,bookmarks=false]{hyperref}
\newcommand{\ie}{i.e.\ }
\newcommand{\cf}{cf.\ }
\newcommand{\resp}{respectively\ }
\newcommand{\Subsection}{\subsection}
\providecommand{\nobreakdash}{\nobreak\hbox{-}}
\setlength{\emergencystretch}{2em}
\multlinegap0pt
\usepackage{tikz}
\usetikzlibrary{cd}
\usepackage{mathtools}
\usepackage{amssymb}
\usepackage{tikz}
\usepackage[shortlabels]{enumitem}
\usepackage{stackengine}
\usepackage{float}
\newtheorem{prop}{Proposition}
\newtheorem{lem}{Lemma}
\usepackage{cleveref}
\crefname{prop}{Proposition}{Propositions}
\crefname{lem}{Lemma}{Lemmas}
\newcommand{\Alg}{\mathrm{Alg}\vphantom{z}}
\newcommand{\Cat}{\mathrm{Cat}\vphantom{t}}
\newcommand{\Mfd}{\mathrm{Mfd}\vphantom{d}}
\newcommand{\Mon}{\mathrm{Mon}_{\infty}}
\newcommand{\PreBord}{\mathrm{Bord}\vphantom{d}}
\newcommand{\PreMfd}{\mathrm{PreMfd}\vphantom{d}}
\DeclareMathOperator{\Vect}{Vect}
\newcommand{\bR}{\mathbb{R}}
\newcommand{\cA}{\mathcal{A}}
\newcommand{\cB}{\mathcal{B}}
\newcommand{\cE}{\mathcal{E}}
\newcommand{\cF}{\mathcal{F}}
\newcommand{\cI}{\mathcal{I}}
\DeclareMathOperator*{\colim}{colim}
\newcommand{\del}{\partial}
\DeclareMathOperator{\op}{op}
\newcommand{\sfA}{\mathsf{A}}
\newcommand{\sfB}{\mathsf{B}}
\newcommand{\sfE}{\mathsf{E}}
\newcommand{\sfI}{\mathsf{I}}
\renewcommand{\to}{\mathchoice{\longrightarrow}{\rightarrow}{}{}}
\title{Structured Flow Categories and Twisted Presheaves}
\author{Alice Hedenlund, Trygve Poppe Oldervoll}
\date{Source: arXiv:2603.29576v1 (CC BY 4.0)}
\begin{document}
\maketitle

\noindent\emph{Source and licence.} Structured Flow Categories and Twisted Presheaves. Version arXiv:2603.29576v1 (CC BY 4.0), distributed by arXiv under the Creative Commons Attribution 4.0 International licence, http://creativecommons.org/licenses/by/4.0/, as recorded in the arXiv licence metadata for this exact version and shown on the abstract page https://arxiv.org/abs/2603.29576v1. Full licence notice: \texttt{licenses/arxiv-2603.29576-CC-BY-4.0.txt}.\par\smallskip

\noindent\textit{Editorial scope.} Counted unit: the article's lem lem:Bord\_e->Bord\_is\_trivial (source0.tex lines 5907--5909). The article's complete original proof of it is delivered with it (source0.tex lines 5911--5985, one \texttt{proof} environment of 75 lines). No mathematics is written, repaired or re-derived: the statement and the proof are the article's own lines, in the article's own notation, and no retained line is substituted or re-flowed.  Retained uncounted as the source-local context the proof needs: lines 906--924; lines 5803--5808; lines 5890--5895.  Nothing was omitted from any retained span, so the package carries no omission record.  No retained line contains a citation, so the wrapper carries no bibliography and no bibliography entry.  The retained lines contain the article's own diagram code, which the wrapper typesets with the standard tikz packages; no image file is delivered and the package declares no asset dependency.  Editorial formatting only, in addition: an AMS \texttt{amsart} preamble that restates the article's own declarations and macros used by the retained lines, namely the environments \texttt{prop}, \texttt{lem} on separate counters, together with the standard packages those lines need.  The retained environments print in this document as Proposition 1, Lemma 1 in order of appearance, whereas the article prints the same items with the numbering of its own full document.  The complete native source of the article as distributed in the arXiv e-print for this exact version is bundled unchanged as \texttt{sources/197527/source0.tex}.
	\begin{prop}\label{thm:iterated_fibration}
		Given a morphism of Cartesian fibrations over $\cB$
		\begin{equation*}
			\begin{tikzcd}
				\cE \arrow[rr, "r"] \arrow[dr, "p"'] & & \cF \arrow[dl, "q"] \\
				& \cB &
			\end{tikzcd}
		\end{equation*}
		such that for each edge $b\in \cB$, the map on fibers $r_{b}\colon \cE_{b} \to \cF_{b}$ is a Cartesian fibration, and such that for each morphism $f\colon b\to c$ in $\cB$, the square
		\begin{equation*}
			\begin{tikzcd}
				\cE_{c} \arrow[r, "f^{*}_{\cE}"] \arrow[d, "r_{c}"]      &   \cE_{b}  \arrow[d, "r_{b}"] \\
				%
				\cF_{c}  \arrow[r, "f^{*}_{\cE}"]               &   \cF_{b}
			\end{tikzcd}
		\end{equation*}
		is a morphism of Cartesian fibrations, then $r$ is a Cartesian fibration.
    If each $\cE_{b}\to \cF_{b}$ is a right fibration, the last condition is automatic, and $r$ is also a right fibration.
	\end{prop}

\begin{lem}\label{lem:A=I}
	There is an equivalence of cartesian fibrations
	\begin{align*}
		\cI^{*}\Alg(\Mfd_{\diamond}^{\mu}) \simeq (\cA\circ \iota)^{*}\Alg(\Mfd_{\diamond}^{\mu}) \,.
	\end{align*}
\end{lem}

\begin{lem}\label{lem:Mfd_e->Mfd_is_right_fib}
	Pushforward by $\pi$ determines a right fibration
	\begin{align*}
		(\cA_{\leq}\circ \iota)^{*}\Alg(\Mfd_{e}) \to (\cA_{\leq}\circ \iota)^{*}\Alg(\Mfd_{\diamond})\,.
	\end{align*}
\end{lem}

\begin{lem}\label{lem:Bord_e->Bord_is_trivial}
	The map $\PreBord_{e} \to \PreBord$ of semi-simplicial spaces obtained from unstraightening the right fibration in \Cref{lem:Mfd_e->Mfd_is_right_fib} is a trivial fibration.
\end{lem}

\begin{proof}
	We need to solve the lifting problem
	\begin{equation*}
		\begin{tikzcd}
			\del \Delta^{n}_{s} \arrow[r] \arrow[d]      &    \PreBord_{e} \arrow[d] \\
			%
			\Delta^{n}_{s}  \arrow[r]     \arrow[ur,dashed]           &     \PreBord \,.
		\end{tikzcd}
	\end{equation*}
	For $n=0$, this is trivial.
  For $n>0$, it suffices as in the proof of horn filling to extend the lift over the minimal object of $\sfI_{[n]}(0,n)$, i.e., to solve the lifting problem
	\begin{equation*}
		\begin{tikzcd}
			\del \sfA_{[n]}(0,n)^{\op} \arrow[r] \arrow[d]      &   \Mfd_{e}  \arrow[d] \\
			%
			\sfA_{[n]}(0,n)^{\op}  \arrow[r]   \arrow[ur,dashed]             &    \Mfd \,.
		\end{tikzcd}
	\end{equation*}
	As in the proof of \Cref{lem:Mfd_e->Mfd_is_right_fib}, the forgetful functor $\Mon \to \Cat_{\infty}$ preserves sifted colimits, so the colimit $\Mfd_{e}$ of $F$ can be computed in $\Cat_{\infty}$ as follows.
  First unstraighten $F$ to a coCartesian fibration $\widetilde{\Mfd}_{e} \to \sfB\Vect$.
  A morphism $N_{0} \to N_{1}$ in $\widetilde{\Mfd}_{e}$ covering the morphism $V$ in $\sfB\Vect$ is precisely a map $N_{0}\oplus V \to N_{1}$ in $\PreMfd_{e}$.
  Then we can obtain the colimit $\Mfd_{e}$ by localizing $\widetilde{\Mfd}_{e}$ at coCartesian edges.
  The coCartesian transport of any $V\in \Vect$ is given by acting with $V$, and since any $V$ is equivalent to some $\bR^{n}$, we have
	\begin{align*}
		\Mfd_{e} \simeq \colim\left(\widetilde{\Mfd}_{e} \xrightarrow{\oplus \bR} \widetilde{\Mfd}_{e} \xrightarrow{\oplus \bR} \dots \right).
	\end{align*}
	Since $\del \sfA_{[n]}(0,n)$ is a finite category, any map to a filtered colimit factors through some finite step.
  The inclusion $\Vect \to \PreMfd_{e}$ is split by the functor which remembers only the vector space $V$.
  We let $\sfE\Vect \to \sfB\Vect$ denote the unstraightening of the action of $\Vect$ on itself.
  We then have an iterated pullback diagram
	\begin{equation*}
		\begin{tikzcd}
			\PreMfd_{e} \arrow[r] \arrow[d]      &   \Vect  \arrow[d]  \arrow[r] &    \ast \arrow[d] \\
			%
			\widetilde{\Mfd}_{e}  \arrow[r]               &   \sfE\Vect  \arrow[r]            &  \sfB \Vect \, ,
		\end{tikzcd}
	\end{equation*}
	so by pasting, the left-hand square is a pullback.
  The functor $\widetilde{\Mfd}_{e}\to \sfE \Vect$ is a morphism of coCartesian fibrations over $\sfB \Vect$, which in the fiber over a point in $\sfB\Vect$ becomes $\PreMfd_{e}\to \Vect$ whose target is a space and therefore a coCartesian fibration.
  Moreover, coCartesian transport along a morphism $V$ in $\sfB\Vect$ induces the square
	\begin{equation*}
		\begin{tikzcd}
			\PreMfd_{e} \arrow[r, "-\oplus V"] \arrow[d]      &   \PreMfd_{e}  \arrow[d] \\
			%
			\Vect  \arrow[r, "-\oplus V"]               &   \Vect
		\end{tikzcd}
	\end{equation*}
	which is a morphism of coCartesian fibrations.
  It then follows from \Cref{thm:iterated_fibration} that $\widetilde{\Mfd}_{e}\to \sfE\Vect$ is a coCartesian fibration.
  Any functor $F\colon \del \sfA_{[n]}(0,n)^{\op} \to \sfE \Vect$ extends to the cone $\sfA_{[n]}(0,n)$ by compatibly embedding all the $F(\sigma)$ in some large inner product space $V$, and using direct sum with the orthogonal complement to construct morphisms.
  Note that we are free to change this lift by including $V$ in some larger $V'$.
  After picking such a lift, it suffices to solve the lifting problem
	\begin{equation*}
		\begin{tikzcd}
			\del \sfA_{[n]}(0,n)^{\op} \arrow[r] \arrow[d]      &   \widetilde{\Mfd}_{e} \arrow[d] \\
			%
			\sfA_{[n]}(0,n)^{\op}  \arrow[r]  \arrow[ur,dashed]              &   \sfE \Vect \times \Mfd_{\diamond} \,.
		\end{tikzcd}
	\end{equation*}
	Since the projection to $\Mfd_{\diamond}$ is invariant under coCartesian transport along morphisms in $\sfE \Vect$, we can use coCartesian transport to the fiber over the terminal object to reduce to the lifting problem
	\begin{equation*}
		\begin{tikzcd}
			\del \sfA_{[n]}^{\op} \arrow[r] \arrow[d]      &   \PreMfd_{e,V}  \arrow[d] \\
			%
			\sfA_{[n]}^{\op}  \arrow[r] \arrow[ur,dashed]              &   \Mfd_{\diamond}
		\end{tikzcd}
	\end{equation*}
	where $\PreMfd_{e,V}$ denotes the fiber of $\PreMfd_{e}\to \Vect$ over $V$.
  This data corresponds to a manifold $M$ stratified by $\sfA_{[n]}(0,n)$, a compatible system of vector bundles $N_{\sigma} \to \del^{\sigma}M$ and maps $N_{\sigma} \to V$, and a compatible system of maps $\del^{\sigma}M\to \del^{\sigma}I^{n}$ where we equip $I^{n}$ with the $\sfA_{[n]}(0,n)$-stratification exhibited by \Cref{lem:A=I}.

	The maps to $\del^{\sigma}I^{n}$ induce a coherent normal framing on the boundary of $M$, which can be extended to all of $M$.
  Using the collar coordinates defined by this coherent normal framing, we can extend the map $\del M \to \del I^{n}$ to a smooth map $M\to I^{n}$ transversely inducing the stratification of $M$.
  After stabilizing to a sufficiently large $V$, we can further extend this to an embedding $M\to I^{n}\times V$.
  The tubular neighborhood $N_{\sigma}$ can similarly be extended to all of $M$ by using collar coordinates near the boundary.
\end{proof}

\end{document}
